% Opracowanie własne; schemat procedury opisanej w rozdziałach 3–5.
\begin{minipage}[t]{0.31\linewidth}\centering\textbf{Trening foldu}\par\medskip
\fbox{\parbox{0.85\linewidth}{\centering $\mathcal{T}_{25\%}\subset\mathcal{T}_{50\%}\subset\mathcal{T}_{100\%}$}}\par\smallskip$\downarrow$\par\smallskip\fbox{\parbox{0.85\linewidth}{\centering Przekształcenia dopasowane osobno na podzbiorze}}\end{minipage}
\hfill
\begin{minipage}[t]{0.31\linewidth}\centering\textbf{Trzy warunki}\par\medskip
\fbox{\parbox{0.85\linewidth}{\centering Pełna regularyzacja}}\par\medskip\fbox{\parbox{0.85\linewidth}{\centering Bez dropoutu}}\par\medskip\fbox{\parbox{0.85\linewidth}{\centering Bez zaniku wag}}\end{minipage}
\hfill
\begin{minipage}[t]{0.31\linewidth}\centering\textbf{Stała ocena}\par\medskip
\fbox{\parbox{0.85\linewidth}{\centering Te same indeksy walidacji}}\par\smallskip$\downarrow$\par\smallskip\fbox{\parbox{0.85\linewidth}{\centering Ten sam test zewnętrzny}}\par\smallskip$\downarrow$\par\smallskip\fbox{\parbox{0.85\linewidth}{\centering Porównania dla sparowanych ziaren}}\end{minipage}\par\medskip Trzy warunki uczenia stosuje się do każdego podzbioru; nie są kolejnymi etapami jednej sieci.
